%% LyX 2.4.3 created this file.  For more info, see https://www.lyx.org/.
%% Do not edit unless you really know what you are doing.
\documentclass[letterpaper,english,reprint, aps]{revtex4-1}
\usepackage[T1]{fontenc}
\usepackage[latin9]{inputenc}
\setcounter{secnumdepth}{3}
\usepackage{amsmath}
\usepackage{amssymb}

\makeatletter

%%%%%%%%%%%%%%%%%%%%%%%%%%%%%% LyX specific LaTeX commands.
\pdfpageheight\paperheight
\pdfpagewidth\paperwidth


\makeatother

\usepackage{babel}
\begin{document}
\title{Sign-problem-free Bosonic Auxiliary Field Quantum Monte Carlo}
\author{Yin-Kai Yu}
\affiliation{University of Chinese Academy of Sciences, Beijing 100049, China}
\author{Zhong-Chao Wei}
\email{zcwei@thp.uni-koeln.de}

\affiliation{Qingdao Binhai Univeristy, Qingdao 266555, China}
\author{Lei Wang}
\email{wanglei@iphy.ac.cn}

\affiliation{Beijing National Laboratory for Condensed Matter Physics and Institute
of Physics,~\\
Chinese Academy of Sciences, Beijing 100190, China}
\affiliation{Songshan Lake Materials Laboratory, Dongguan, Guangdong 523808, China}
\author{Zi-Xiang Li}
\email{zixiangli@iphy.ac.cn}

\affiliation{Beijing National Laboratory for Condensed Matter Physics and Institute
of Physics,~\\
Chinese Academy of Sciences, Beijing 100190, China}
\affiliation{University of Chinese Academy of Sciences, Beijing 100049, China}
\date{\today}
\begin{abstract}
This is a note on reflection positivity, time-reversal symmetry and
related sign-problem-free models in bosonic auxiliary field quantum
Monte Carlo(BAFQMC).
\end{abstract}
\keywords{sign problem, bosonic system, quantum Monte Carlo}
\maketitle

\section{Introduction}

The concept of reflection positivity was originally introduced in
the context of axiomatic quantum field theory\citep{osterwalder_axioms_1973}.
It has then become a standard mathematical tool in statistical physics\citep{kondratiev_reflection_2006,biskup_reflection_2009}.
Its main application in strongly correlated systems started from Lieb\textquoteright s
work on the spin reflection positivity in the Hubbard model\citep{lieb_two_1989}.
Majorana reflection positivity was proposed\citep{jaffe_reflection_2015,jaffe_characterization_2016}
and then applied to interacting spinless fermion models\citep{wei_ground_2015}.
The connection between Majorana reflection positivity and sign-problem-free
quantum lattice fermion models was then discovered in Ref.~\citep{wei_majorana_2016}.

In addition to this, the time-reversal symmetry of fermion models
also implies the absence of the sign problem.

In this paper, we explore the sign problem of quantum lattice boson
models from the perspective of reflection positivity and time-reversal
invariance. 

\section{The trace formula in BAFQMC}

In BAFQMC algorithms for interacting quantum lattice boson models,
interaction terms are decoupled by auxiliary fields into bosonic quadratic
forms. After integrating out the bosonic degrees of freedom, an action
in terms of auxiliary fields can be obtained. This action can be treated
with random sampling numerically. The sampling weight for each configuration
of auxiliary fields has the form\citep{thomas_character_2008}
\begin{align}
p & =\mathrm{tr}\left(e^{h_{1}}e^{h_{2}}\cdots e^{h_{k}}\right)\nonumber \\
 & =\det\left(I-e^{A_{1}}e^{A_{2}}\cdots e^{A_{k}}\right)^{-1/2}.\label{eq:TraceFormula}
\end{align}

The sampling weight expressed in Eq.~(\ref{eq:TraceFormula}) is
the main object of study in this work. Here $h_{i}=-\frac{1}{2}\chi^{T}\Omega A_{i}\chi$
($i=1,2,\dots k$) denote a set of bosonic quadratic forms. They come
from both single-body terms and decoupling of interaction terms using
auxiliary fields, and vary as the configuration of auxiliary fields
changes. We have employed the basis $\chi=\left(\begin{array}{c}
a_{1}\\
a_{2}\\
\vdots\\
a_{N}\\
a_{1}^{+}\\
a_{2}^{+}\\
\vdots\\
a_{N}^{+}
\end{array}\right)$, with canonical commutation relations of bosonic creation and annihilation
operators $\left[a_{i},\ a_{j}\right]=0$, $\left[a_{i}^{+},\ a_{j}^{+}\right]=0$,
$\left[a_{i},\ a_{j}^{+}\right]=\delta_{ij}$($i,j=1,2,\dots N$),
and $\Omega=\left(\begin{array}{cc}
0 & I_{N}\\
-I_{N} & 0
\end{array}\right)$. In principle $h_{i}$ might contain an arbitrary particle number
conserving part and an arbitrary pairing part unless we specify a
particular example. That is to say the coefficients $A_{i}$ could
be arbitrary complex square matrices that satisfies $\Omega A+A^{T}\Omega=0$.

If there is no paring part in the expression of each $h_{i}$, let
$\chi_{0}=\left(\begin{array}{c}
a_{1}\\
a_{2}\\
\vdots\\
a_{N}
\end{array}\right)$, $h_{i}=\chi_{0}^{\dagger}B_{i}\chi_{0}$, then Eq.~(\ref{eq:TraceFormula})
degenerates to
\begin{align}
p & =\mathrm{tr}\left(e^{h_{1}}e^{h_{2}}\cdots e^{h_{k}}\right)\nonumber \\
 & =\det\left(I-e^{B_{1}}e^{B_{2}}\cdots e^{B_{k}}\right)^{-1}.\label{eq:TraceFormula2}
\end{align}

Eq.~(\ref{eq:TraceFormula2}) does not only holds for ``physical''
circumstances when the moduli of the eigenvalues of $e^{B_{1}}e^{B_{2}}\cdots e^{B_{k}}$
are less than one, but also holds in the sense of regularization for
``unphysical'' circumstances. See the Supplemental Material for
Gaussian cut-off renormalization of the sampling weights. In the same
vein, Eq.~(\ref{eq:TraceFormula}) should also be interpreted as
the regularized sampling weight. In BAFQMC simulations, when the chemical
potential takes unphysical values, the partition function is usually
a regularized finite number instead of infinity, and the expectation
values of observables also take regularized values.

If no restriction is put on $h_{i}$, the bosonic Gaussian integral
$p$ could be non-positive, due to both the complex nature of the
coefficients and the two-valuedness of the metapletic representation.
Statistical sampling methods may fail to obtain desired physical quantities
with useful accuracy at reasonable cost under those circumstances.
This is the so-called sign problem in BAFQMC methods.

\section{Calculation of expectation values of observables}

Let $\left\{ X_{s}\right\} $($s=1,2,\dots,2n$) be a set of bosonic
creation or annihilation operators, and $Z=\mathrm{tr}\left(e^{h_{1}}e^{h_{2}}\cdots e^{h_{k}}\right)$.
Without losing generality, we may assume $Z\neq\pm\infty$. Define
the Green function as
\begin{equation}
G_{ij}=\frac{1}{Z}\mathrm{tr}\left\{ \left[\theta\left(i-j\right)X_{i}X_{j}+\theta\left(j-i\right)X_{j}X_{i}\right]e^{h_{1}}e^{h_{2}}\cdots e^{h_{k}}\right\} ,
\end{equation}
where $\theta$ is the step function. And define a symmetric matrix(with
zero diagonal elements) $\tilde{G}$ by $\tilde{G}_{ij}=G_{ij}$ when
$i>j$. Then by Wick's theorem we have
\begin{equation}
\frac{1}{Z}\mathrm{tr}\left\{ X_{2n}\cdots X_{2}X_{1}e^{h_{1}}e^{h_{2}}\cdots e^{h_{k}}\right\} =\frac{\text{1}}{n!2^{n}}\sum_{\sigma\in S_{2n}}\prod_{i=1}^{n}\tilde{G}_{\sigma\left(2i-1\right)\sigma\left(2i\right)},
\end{equation}
where $S_{2n}$ is the symmetric group.

For the calculation of the Green function, we have 
\begin{equation}
\textrm{tr}\left(\chi_{\alpha}\chi_{\beta}e^{h_{1}}e^{h_{2}}\cdots e^{h_{k}}\right)/Z=\left(\frac{1}{I-e^{A_{1}}e^{A_{2}}\cdots e^{A_{k}}}\Omega\right)_{\alpha\beta}.
\end{equation}

If there is no paring part, these calculations can be further simplified
and have already been discussed by S.W. Zhang.

\section{Reflection positivity and sign-problem-free BAFQMC simulations }

Let $b_{i}$, $b_{i}^{+}$($i=1,2,\dots N$) and $c_{i}$, $c_{i}^{+}$($i=1,2,\dots N$)
be two sets of bosonic annihilation/creation operators. Then any operator
$O$ can be expressed by
\begin{equation}
O=\sum_{\alpha,\beta}W_{\alpha\beta}\Gamma_{\alpha}^{\left(1\right)}\Gamma_{\beta}^{\left(2\right)},
\end{equation}
where $\Gamma_{\alpha}^{\left(1\right)}=b_{i_{1}}^{+}b_{i_{2}}^{+}\cdots b_{i_{m}}^{+}b_{j_{1}}b_{j_{2}}\cdots b_{j_{n}}$,
$\text{\ensuremath{\alpha=\left\{ \left\{ i_{1},i_{2},\dots,i_{m}\right\} \left\{ j_{1},j_{2},\dots,j_{n}\right\} \right\} }}$,
$\Gamma_{\beta}^{\left(2\right)}=c_{i_{1}}^{+}c_{i_{2}}^{+}\cdots c_{i_{s}}^{+}c_{j_{1}}c_{j_{2}}\cdots c_{j_{t}}$,
$\text{\ensuremath{\beta=\left\{ \left\{ i_{1},i_{2},\dots,i_{s}\right\} \left\{ j_{1},j_{2},\dots,j_{t}\right\} \right\} }}$.
Each index $i_{k}$ or $j_{l}$ can repeat zero or several times in
this definition of $\alpha$ and $\beta$.

Consider the reflection operation $R$, which is an anti-linear transformation
that sends $c\mapsto d$, $d\mapsto c$ and $i\mapsto-i$. We say
an operator $O$ is reflection symmetric if and only if it does not
change under $R$, namely $W=W^{\dagger}$. We say a reflection symmetric
operator $O$ is reflection positive if and only if $W\geq0$. Clearly
the set of reflection positive operators is closed under operator
summation and multiplication, and any reflection positive operator
has non-negative trace.

To relate reflection positivity to sign-problem-free models, consider
the following type of bosonic quadratic forms:
\begin{align}
h & =\left(\begin{array}{cccc}
b^{T} & b^{+T} & c^{T} & c^{+T}\end{array}\right)\left(\begin{array}{cc}
D & F\\
F^{T} & \bar{D}
\end{array}\right)\left(\begin{array}{c}
b\\
b^{+}\\
c\\
c^{+}
\end{array}\right),
\end{align}
where $F\geq0$ is a positive hermitian matrix. Then $e^{h}$ is reflection
positive. If all the $e^{h_{i}}$ in Eq.~(\ref{eq:TraceFormula})
take this form, then we will have $p\geq0$ and thus there is no sign
problem in this BAFQMC simulation.

However, this argument holds only for bare sampling weights. For regularized
sampling weights, $e^{A_{1}}e^{A_{2}}\cdots e^{A_{k}}$ in Eq.~(\ref{eq:TraceFormula})
can be decomposed into a direct sum of an element in a contraction
semigroup and an element in an expansion semigroup. Therefore we still
have $p\geq0$\citep{wei_semigroup_2024}. See the Supplemental Material
for details.

\section{Time-reversal invariance and sign-problem-free BAFQMC simulations}

Similar to the reflection operation, the time-reversal operation $T$
for bosons is an anti-linear transformation that sends $c\mapsto id$,
$d\mapsto-ic$ and $i\mapsto-i$. We consider the time-reversal symmetric
bosonic quadratic forms without pairing terms:
\begin{equation}
h=\left(\begin{array}{cc}
b^{+T} & c^{+T}\end{array}\right)\left(\begin{array}{cc}
D & F\\
-\bar{F} & \bar{D}
\end{array}\right)\left(\begin{array}{c}
b\\
c
\end{array}\right),
\end{equation}
where $D$ and $F$ are arbitrary complex square matrices. If all
the $h_{i}$ in Eq.~(\ref{eq:TraceFormula2}) take this form, then
all the eigenvalues of $e^{B_{1}}e^{B_{2}}\cdots e^{B_{k}}$ will
appear in complex conjugate pairs due to Kramers degeneracy, so we
will have $p\geq0$ and thus there is no sign problem in this BAFQMC
simulation.

\section{Simulation of sign-problem-free models}

Consider a quantum lattice boson model with single body term 

\begin{equation}
H_{K}=b^{+T}Bb+c^{+T}\bar{B}c,
\end{equation}
and interaction term
\begin{equation}
H_{U}=U_{1}\sum_{i}\left(b_{i}^{+}b_{i}+c_{i}^{+}c_{i}\right)^{2}+U_{2}\sum_{i}\left(b_{i}^{+}b_{i}-c_{i}^{+}c_{i}\right)^{2},
\end{equation}
where $B=B^{\dagger}$, $U_{1}\leq0$, $U_{2}\geq0$. This model can
be seen as an example of both reflection positivity and time-reversal
invariance.
\begin{acknowledgments}
We would like to thank...
\end{acknowledgments}

\bibliographystyle{apsrev4-1}
\bibliography{RPBAFQMC}
\newpage{}

\appendix

\section*{Supplemental Material}

\section{Gaussian cut-off renormalization of the sampling weights}

The series to regularize is $p=\sum_{n_{1}=0}^{\infty}\sum_{n_{2}=0}^{\infty}\dots\sum_{n_{N}=0}^{\infty}r_{1}^{n_{1}}r_{2}^{n_{2}}\dots r_{N}^{n_{N}}=\mathrm{tr}\left(e^{h_{1}}e^{h_{2}}\cdots e^{h_{k}}\right)$,
where $h_{i}=\chi_{0}^{\dagger}B_{i}\chi_{0}$, $\chi_{0}=\left(\begin{array}{c}
a_{1}\\
a_{2}\\
\vdots\\
a_{N}
\end{array}\right)$, $r_{1},r_{2},\dots,r_{N}$ are the $N$ eigenvalues of $e^{B_{1}}e^{B_{2}}\cdots e^{B_{k}}$,
and $n_{1},n_{2},\dots,n_{N}$ are boson particle numbers. Let $a_{n_{1}n_{2}\dots n_{N}}=r_{1}^{n_{1}}r_{2}^{n_{2}}\dots r_{N}^{n_{N}}$,
$g\left(a_{n_{1}n_{2}\dots n_{N}}\right)=\exp\left(-\left|a_{n_{1}n_{2}\dots n_{N}}\right|^{2}/\Lambda^{2}\right)$.
Using Laurent\textquoteright s expansion with respect to $\Lambda$,
we have
\begin{align*}
 & p_{\Lambda}\left(\Lambda,r_{1},r_{2},\dots,r_{N}\right)\\
 & =\sum_{n_{1}=0}^{\infty}\sum_{n_{2}=0}^{\infty}\dots\sum_{n_{N}=0}^{\infty}a_{n_{1}n_{2}\dots n_{N}}g\left(a_{n_{1}n_{2}\dots n_{N}}\right)\\
 & =P\left(\Lambda,r_{1},r_{2},\dots,r_{N}\right)+\prod_{i=1}^{N}\frac{1}{1-r_{i}}+O\left(\frac{1}{\Lambda}\right).
\end{align*}
As $\Lambda\rightarrow\infty$, the terms in $P\left(\Lambda,r_{1},r_{2},\dots,r_{N}\right)$
either vanish or diverge no slower than $O\left(\Lambda\right)$.
Define

\begin{align*}
 & p_{reg}\left(r_{1},r_{2},\dots,r_{N}\right)\\
 & =\prod_{i=1}^{N}\frac{1}{1-r_{i}}\\
 & =\det\left(I-e^{B_{1}}e^{B_{2}}\cdots e^{B_{k}}\right)^{-1}.
\end{align*}

In the calculation of partition function, $Z=\sum_{\left\{ \phi_{i}\right\} }p_{\left\{ \phi_{i}\right\} }$,
where $\left\{ \phi_{i}\right\} $ stands for configurations of auxiliary
fields. Let's calculate $Z_{\Lambda}=\sum_{\left\{ \phi_{i}\right\} }p_{\Lambda\left\{ \phi_{i}\right\} }$
instead and then let $\Lambda\rightarrow\infty$, $Z=\lim_{\Lambda\rightarrow\infty}Z_{\Lambda}$.
For physical(low enough) choice of the chemical potential, the summation
of $P\left(\Lambda,r_{1},r_{2},\dots,r_{N}\right)$ must lead to a
exact cancellation, otherwise $Z_{\Lambda}$ will diverge as $\Lambda\rightarrow\infty$.
So we have $Z=\sum_{\left\{ \phi_{i}\right\} }p_{reg\left\{ \phi_{i}\right\} }$.
For unphysical choice of the chemical potential, $Z$ is infinite
but $\sum_{\left\{ \phi_{i}\right\} }p_{reg\left\{ \phi_{i}\right\} }$is
still finite(unless the chemical potential takes some critical value).

\section{Positivity of regularized sampling weights}

For any bosonic quadratic form $h=-\frac{1}{2}\chi^{T}\Omega A\chi$,
assume $A=\left(\begin{array}{cccc}
B_{11} & B_{12} & C_{11} & C_{12}\\
B_{21} & B_{22} & C_{12}^{T} & C_{22}\\
D_{11} & D_{12} & -B_{11}^{T} & -B_{21}^{T}\\
D_{12}^{T} & D_{22} & -B_{12}^{T} & -B_{22}^{T}
\end{array}\right),$ where $C_{11}=C_{11}^{T}$, $C_{22}=C_{22}^{T}$, $D_{11}=D_{11}^{T}$,
$D_{22}=D_{22}^{T}$. Let us introduce two $Z_{2}$ anti-linear symmetric
transformations, expressed in terms of $J_{1}=\left(\begin{array}{cccc}
0 & I_{N/2} & 0 & 0\\
I_{N/2} & 0 & 0 & 0\\
0 & 0 & 0 & I_{N/2}\\
0 & 0 & I_{N/2} & 0
\end{array}\right)$, $J_{2}=\left(\begin{array}{cccc}
0 & I_{N/2} & 0 & 0\\
-I_{N/2} & 0 & 0 & 0\\
0 & 0 & 0 & I_{N/2}\\
0 & 0 & -I_{N/2} & 0
\end{array}\right)$. Reflection symmetry requires $J_{1}^{T}AJ_{1}=\bar{A}$. Assume
that $J_{2}\Omega$ plays the role of the metric of the contraction
semigroup and the expansion semigroup. For the generators of the maximal
subgroups of the semigroups, we have $J_{2}^{T}AJ_{2}=\bar{A}$, so
$A=\left(\begin{array}{cccc}
B_{11} & 0 & C_{11} & 0\\
0 & \bar{B_{11}} & 0 & \bar{C}_{11}\\
D_{11} & 0 & -\bar{B}_{11} & 0\\
0 & \bar{D}_{11} & 0 & -B_{11}
\end{array}\right)$, where $C_{11}=C_{11}^{T}$, $D_{11}=D_{11}^{T}$. For the generators
in the invariant cones of the semigroups, we have we have $J_{2}^{T}AJ_{2}=-\bar{A}$,
so $A=\left(\begin{array}{cccc}
0 & B_{12} & 0 & C_{12}\\
\bar{B}_{12} & 0 & \bar{C}_{12} & 0\\
0 & D_{12} & 0 & -B_{12}^{\dagger}\\
\bar{D}_{12} & 0 & -B_{12}^{T} & 0
\end{array}\right),$where $C_{12}=C_{12}^{\dagger}$, $D_{12}=D_{12}^{\dagger}$, and
$J_{2}\Omega A=\left(\begin{array}{cccc}
\bar{D}_{12} & 0 & -B_{12}^{T} & 0\\
0 & -D_{12} & 0 & B_{12}^{\dagger}\\
-\bar{B}_{12} & 0 & -\bar{C}_{12} & 0\\
0 & B_{12} & 0 & C_{12}
\end{array}\right)$. The reflection positivity condition is equivalent to $\left(\begin{array}{cc}
C_{12} & B_{12}\\
B_{12}^{\dagger} & -D_{12}
\end{array}\right)\geq0$, which means $e^{A}$ can be decomposed into a direct sum of an element
in the contraction semigroup and an element in the expansion semigroup.
\end{document}
